\documentclass[11pt,a4paper]{article}
\usepackage{fontspec}
\setmainfont{Noto Serif}
\setsansfont{Noto Sans}
\usepackage{babel}
\babelprovide[main,import]{russian}
\usepackage{amsmath,amssymb,mathtools,booktabs,geometry,fancyhdr,microtype,xcolor,hyperref}
\geometry{left=21mm,right=21mm,top=21mm,bottom=21mm}
\definecolor{ink}{HTML}{173D57}
\hypersetup{colorlinks=true,urlcolor=ink,linkcolor=ink,pdftitle={BGS-F. Возвраты и спектральный перенос}}
\pagestyle{fancy}\fancyhf{}\lhead{\small BGS-F}\rhead{\small Возвраты и спектральный перенос}\cfoot{\thepage}
\setlength{\headheight}{15pt}\setlength{\parindent}{1em}\setlength{\parskip}{2.5pt}
\setlength{\emergencystretch}{3em}\tolerance=3000
\newcommand{\Tr}{\operatorname{Tr}}
\newcommand{\EE}{\mathbb E}
\newcommand{\R}{\mathbb R}
\newcommand{\Z}{\mathbb Z}
\newcommand{\C}{\mathbb C}
\newcommand{\si}{\mathrm{Sine}_2}
\newcommand{\eps}{\varepsilon}
\newcommand{\sinc}{\operatorname{sinc}}
\newcommand{\sect}[1]{\section*{\sffamily\color{ink}#1}}
\newcommand{\subt}[1]{\par\medskip\noindent\textbf{#1.}\quad}
\newcommand{\T}{\mathbb T}
\newcommand{\dist}{\operatorname{dist}}
\newcommand{\Ad}{\operatorname{Ad}}
\begin{document}
\noindent{\small\sffamily\color{ink} МАТЕМАТИЧЕСКАЯ ФИЗИКА\hfill 20 НОЯБРЯ 2026}
\begin{center}{\LARGE\sffamily Гейзенберговские возвраты до $7^{59}$}\\[4pt]
{\large Второй порядок и перенос уровней между узлами}
\end{center}
Рассматриваем исходные $A_N,M_N$, $N=7^k$, и $U_N(t)=e^{itA_N}M_N$. Физической силе $u_N$ соответствует $t=Nu_N$. В основной задаче $t=N/10^{80}$. Время скалярного возврата линейной карты равно $m=3N/7$.

\textbf{Исходный фиксированный кик.} Для всех $2\le k\le59$ при $t=N/10^{80}$ и $m=3N/7$
\[
                 |\Tr U_N(t)^m|>0.957N,\qquad
                         \Delta_{2,N}(b_*)>0.053 .
\]
Один ограниченный непрерывный локальный функционал из предыдущей части отличается по ожиданию от $\si$ более чем на $0.76$ на всём этом диапазоне. Старшая доказанная размерность $7^{59}>7\cdot10^{49}$. Для $k=57,58$ нормированный след превышает $0.99999$. Границы получены для неизменной силы. На составной размерности $N>3\cdot10^{69}$ при времени $m>4\cdot10^{15}$ доказано $|\Tr U_N^m|>0.9915N$ и $\Delta_{2,N}(b_*)>0.582$. Доля времени в полной размерности в этом примере меньше $3/7$.

\textbf{Вспомогательные убывающие силы.} На той же последовательности размерностей доказано:
\begin{itemize}
 \item при $u_N=o(N^{-8/5})$ нормированный обычный след на времени $3N/7$ по модулю стремится к 1;
 \item при $u_N=o(N^{-22/13})$ весь локальный процесс имеет арифметический предел. Это случайно сдвинутая решётка шага $7/3$ с кратностями $2(1+\nu_7(B+j))$, где $B$ равномерен в $\Z_7$;
 \item при $u_N=o(N^{-7/4}(\log N)^{-1/20})$ полный парный тест сходится к той же положительной константе, что у невозмущённой карты. Её заключающий интервал $(0.08443,0.08656)$ вычислен в предыдущей части.
\end{itemize}
Все три области шире ранее доказанных для соответствующих утверждений. Для этого подсчитывается число совпадений сумм двух вейлевых меток в коротких временных блоках. Затем две интеграции по частям дают оценку вращения полного произведения киков, а спектральный перенос контролируется числом уровней, пересёкших середину промежутка между узлами.

\textbf{Статус.} Для неизменного $\tau$ установлена новая конечная полоса. Сохраняется ли скалярный гейзенберговский возврат на бесконечном ряду $7^k$ при той же силе, остаётся открытым.
\newpage
\sect{1. Координаты орбиты и нормы первого порядка}
Пусть $p=7$, $N=p^k$, $m=3N/7$, $\|X\|_{\rm n}^2=N^{-1}\Tr X^*X$. Пишем $A=\sum_d c_d T(d)$, где $T(d)$ --- модульные операторы Вейля, ортонормированные в этой норме, $Q(d)=d_1^2-2d_2^2$. На ненулевом носителе $A$ две координаты $d$ не имеют общего множителя 7, кроме меток $(0,\pm7)$. Для каждой метки $d=p^\alpha d'$ примитивная часть имеет $Q(d')$, обратимое по модулю 7. Это проверяется непосредственно из 32 коэффициентов исходного символа. Кроме того,
\[
              \sum_d |c_d|\le 593/500<6/5 .
\]
В кольце $\Z/p^k\Z$ берём корень $\rho^2=2$, $\rho\equiv3\bmod7$. В расщеплённых координатах $(x,y)=(d_1+\rho d_2,d_1-\rho d_2)$ преобразование $S$ умножает $x$ на $\gamma=3+2\rho$, а $y$ на $\gamma^{-1}$. По модулю $p^b$ элемент $\gamma$ имеет порядок $3p^{b-1}$ для всякого $1\le b\le k$. Это следует из $S^3-I=14\bigl(\begin{smallmatrix}7&10\\5&7\end{smallmatrix}\bigr)$ и биномиального поднятия. Порядок нечётный, поэтому $\gamma^2$ имеет тот же порядок.

Для $A_j=M^j A M^{-j}=\sum_d c_d T(S^jd)$ определим $S_l=\sum_{j=j_0}^{j_0+l-1} A_j$. Если $l\le3N/49$, в орбите каждой фиксированной метки все $l$ элементов различны. Если две исходные метки имеют разные целые $Q$, их орбиты не пересекаются при $k\ge3$, поскольку разность норм по модулю $N$ имеет абсолютное значение меньше $N$. Метки $(0,7)$ и $(0,-7)$ лежат в разных орбитах, хотя их $Q$ одинаковы: переход к минусу потребовал бы элемент порядка 2 в циклической группе нечётного порядка. Группируя слагаемые по $Q$ и разделяя эти две метки, получаем
\[
       \|S_l\|_{\rm n}^2
          \le \frac{847747}{2000000}\,l
          < \frac49\,l,\qquad l\le 3N/49 .                          \tag{1}
\]
При вычислении рационального числа для каждого класса остальных меток используется квадрат суммы $|\Re c_d|+|\Im c_d|$. Отдельный вклад двух меток с $Q=-98$ равен $2(9/20)^2$. Старт $j_0$ произволен, поскольку норма инвариантна при сопряжении через $M$.

Для длин до $m$ можно разделить интервал на не более семи частей длины $3N/49$. Из (1) и Коши получаем $\|S_l\|_{\rm n}\le(2/3)\sqrt{7l}$ для любого $l\le m$. В следующем разделе отдельно оцениваются суммы всех произведений второго порядка, где существенен подсчёт совпадающих выходных меток.
\newpage
\sect{2. Сколько пар орбитальных меток имеют одну сумму}
Фиксируем две метки $d,e$ носителя, времена $i,j$ из одного целого интервала длины $l\le m$. Пусть $\alpha=\min(\nu_7(d),\nu_7(e))$, $\delta=|\nu_7(d)-\nu_7(e)|$ и $n=p^K=N/p^\alpha$. Можно поменять местами две метки так, чтобы после деления обеих на $p^\alpha$ первая стала примитивной. Расщеплённые координаты первой орбитальной метки обозначим $(x,y)$, второй --- $(x',y')$. Обе координаты первой обратимы, а координаты второй имеют порядок $\delta\in\{0,1\}$. Нормы $c=xy$, $d_0=x'y'$ постоянны вдоль орбит. Обозначим
\[
               (X,Y)=(x+x',y+y'),\qquad D_0=xy'-x'y .
\]
Кратность представления фиксированного $(X,Y)$ ограничим через $v=\nu_p(D_0)$, усечённую на $K$.

Если $(X,Y)=0\bmod n$, то $v=K$. Иначе пусть $r=\min(\nu_p X,\nu_p Y)<K$. Если $\nu_p Y=r$, координата $z$ первой точки любого другого представления удовлетворяет
\[
        Y z^2-(XY+c-d_0)z+cX=0\bmod p^K.                            \tag{2}
\]
Действительно, первая точка имеет вторую координату $c/z$, а норма второй равна $d_0$. Все коэффициенты (2) делятся на $p^r$ благодаря равенству $XY+c-d_0=Xy+Yx$ в исходном представлении. После деления старший коэффициент обратим, а производная при $z=x$ равна $D_0/p^r$ по модулю $p^{K-r}$. Если $\nu_p X=r$, можно поменять координаты и получить тот же результат с противоположным знаком производной.

Дополняя квадрат, получаем условие $z_1^2=z_0^2$ по модулю $p^{K-r}$ с усечённым порядком $\nu_p z_0=v-r$. Если $2(v-r)<K-r$, решение сравнимо с одним из $\pm z_0$ по модулю $p^{K-v}$ по факторизации разности квадратов и нечётности $p$. Если $2(v-r)\ge K-r$, решения делятся на $p^{\lceil(K-r)/2\rceil}$, а значит, задают один класс по модулю $p^{K-v}$. Аффинная замена с обратимым коэффициентом сохраняет число классов. Итак, выбранная координата первой точки лежит не более чем в двух классах по модулю $p^{K-v}$. В случае $v=K$ это тривиально.

Проекция первой точки пробегает каждый допустимый класс по модулю $p^b$ ровно один раз за период $3p^{b-1}$. Значит, число возможных первых времён не больше
\[
                        2+\frac{14}3\,\frac{l}{n}\,p^v.             \tag{3}
\]
При $v=K$ эта формула покрывает тривиальную границу $l$. Для каждого первого времени вторая точка однозначно определяется, поскольку известна сумма. Поэтому (3) ограничивает кратность пары, если $l\le3N/49$. При $l\le m$ её следует умножить не более чем на 7, по длине наименьшей орбиты. Ограничение на порядок времён может только уменьшить кратность.
\newpage
\sect{3. Сумма квадратов кратностей на коротких блоках}
При фиксированном первом времени и меняющемся втором определитель $D_0$ имеет вид $p^\delta(a\gamma^{-j}-b\gamma^j)$ с обратимыми $a,b$. При $s>\delta$ условие $\nu_p(D_0)\ge s$ задаёт не более одного класса по периоду $3p^{s-\delta-1}$. Действительно, оно равносильно условию на $\gamma^{2j}$ по модулю $p^{s-\delta}$, а квадрат имеет тот же порядок, что и $\gamma$. Следовательно, число подходящих $j$ не больше
\[
                1+\frac l3 p^{1+\delta-s}\qquad(s>\delta).          \tag{4}
\]
Возьмём $k\ge3$, $l\le N/(1000k)$. Пусть $L_w$ есть число \emph{упорядоченных} по времени пар с суммой $w$, например $i\le j$ (после возможной перестановки меток порядок меняется). Требуется оценить $\sum_w L_w^2$. Это сумма величин $L_{w(i,j)}$ по самим парам с тем же ограничением.

Если $l\le3$, то $L_w\le l$ и $\sum_w L_w\le l(l+1)/2$, откуда достаточно границы $2 l^2$. Пусть $l\ge4$. Положим $s_0\ge\delta$ наибольшим допустимым показателем, для которого период при условии (4) не превосходит $l$. При отсутствии такого показателя $>\delta$ берём $s_0=\delta$. Имеем $s_0<K$ и $p^{s_0}\le (p^{1+\delta}/3) l$. Тогда при каждом $i$ не более одного $j$ имеет $v>s_0$. Его вклад в сумму кратностей не больше $l$. Для остальных суммируем (3), используя
\[
 \sum_{\substack{j\\v\le s_0}}p^v
     \le p^\delta l+\sum_{s=\delta+1}^{s_0}
               (p^s-p^{s-1})\left(1+\frac l3p^{1+\delta-s}\right)
     \le p^\delta l(2k+10/3).
\]
Базовая двойка из (3) даёт не больше $l(l+1)$ для упорядоченных пар, исключительные пары дают ещё не больше $l^2$. Так как $p^{\alpha+\delta}\le7$, получаем
\[
    \sum_w L_w^2 \le 2l^2+l+
          \frac{98}{3}\,(2k+10/3)\,\frac{l^3}{N}
                         \le \frac{12}5\,l^2 .                    \tag{5}
\]
Последнее верно при $k\ge3$, $l\ge4$, $l/N\le1/(1000k)$: коэффициент последнего члена после деления на $l^2$ меньше $0.102$.

Для произвольного $l\le m$ используем (3) с множителем 7 и суммируем (4) без отсечения по $s_0$. Тогда $\sum_j p^v\le n+(2K+1)p^\delta l$ для каждого $i$, откуда
\[
                   \sum_w L_w^2 \le C l^2(1+kl/N).                \tag{6}
\]
Здесь $C$ --- абсолютная константа. Этот вариант подходит для блоков, длина которых превышает $N/(1000k)$. Порядок $i\le j$ можно использовать и в нём.
\newpage
\sect{4. Второй порядок произведения унитарных киков}
Все произведения $T(v)T(w)$ имеют метку $v+w$ и фазу единичного модуля. Разные суммы дают ортогональные слагаемые в норме Гильберта--Шмидта. Совпадения сумм оцениваются по модулю. Поэтому для любых $|\lambda_{ij}|\le1$ получаем
\[
   \left\|\sum_{0\le i\le j<l}\lambda_{ij} A_{j_0+j}A_{j_0+i}\right\|_{\rm n}
       \le \frac94\,l\qquad(l\le N/(1000k)).
                                                                     \tag{7}
\]
Сначала применяем (5) к каждой паре исходных меток, затем треугольное неравенство по меткам. Константа $9/4$ больше $(6/5)^2\sqrt{12/5}$. Для $l\le m$ имеем также $C l\sqrt{1+kl/N}$ по (6). Эти оценки допускают дробные коэффициенты с последнего неполного шага.

Чтобы применить (7) ко всему произведению, положим $A(s)=A_{j_0+\lfloor s\rfloor}$ на блоке, $S(s)=\int_0^s A(v)dv$, $B(s)=\int_0^s A(v)S(v)dv$. Пусть $V'(s)=it V(s)A(s)$ и $V(0)=I$. При целом $l$ оператор $V(l)$ есть точное произведение экспонент, унитарное при вещественном $t$. Дважды интегрируя по частям, имеем
\[
 V(l)-I = it V(l)S(l) +t^2 V(l)B(l)
                          -it^3\int_0^l V(s)A(s)B(s)\,ds .         \tag{8}
\]
Функции абсолютно непрерывны, поэтому скачки генератора на границах шагов не создают краевых членов. В $B(s)$ встречаются пары $i\le j$ с весом не больше 1, включая диагональные пары с весом не больше $1/2$. Поэтому $\|B(s)\|_{\rm n}\le(9/4)\lceil s\rceil$ на коротком блоке. Последний интеграл ограничен по норме, используя $\|A\|_{\rm op}\le 6/5$. Из (1),(8)
\[
 \|V(l)-I\|_{\rm n}
       \le \frac23|t|\sqrt l+\frac94 t^2l
                                  +\frac{27}{20}|t|^3l(l+1).
\]

Разбиваем $m$ шагов на полные блоки длины $L$ и возможный остаток $<L$. Ошибки складываются по норме, так как все множители унитарны. Сумма квадратных корней длин не превосходит $m/\sqrt L+\sqrt L$. Для $L\le N/(1000k)$ получаем
\[
 \boxed{\displaystyle
 D_N(t):=\|U_N(t)^m-M_N^m\|_{\rm n}
   \le \frac23|t|\left(\frac m{\sqrt L}+\sqrt L\right)
         +\frac94 t^2 m+\frac{27}{20}|t|^3 m(L+1).}                 \tag{9}
\]
Для любого $1\le L\le m$ справедлива более грубая граница
\[
  D_N(t)\le C\left[
     |t|\left(\frac m{\sqrt L}+\sqrt L\right)
     +t^2m\sqrt{1+kL/N}
     +|t|^3m(L+1)\sqrt{1+kL/N}\right] .                              \tag{10}
\]
Здесь в первом порядке применена длинная версия (1), а во втором --- (6). Обе формулы верны на исходном $A_N$.
\newpage
\sect{5. Исходная сила и доказанный диапазон}
Положим $t=N/10^{80}$ и $m=3N/7$. Для $k=57,58,59$ соответственно выберем квадратные блоки $L=x^2$ с
\[
                 x=10^{21},\qquad 3\cdot10^{20},\qquad
                                      78\cdot10^{18}.
\]
Во всех трёх случаях $L\le N/(1000k)$, как проверяется целыми числами. Поэтому правая часть (9) целиком рациональна. Получены верхние границы
\[
\begin{array}{c|ccc}
 k&57&58&59\\ \hline
 D_N(t)\ \text{меньше}&10^{-5}&0.0017&0.291.
\end{array}
\]
Остальные размерности $2\le k\le56$ покрываются прежней границей $D_N(t)<0.204$ из части об обычном следе. Пусть $M_N^m=\mu_N I$. Тогда
\[
          \Re\!\left(\overline{\mu_N}\frac{\Tr U_N(t)^m}N\right)
                              =1-\frac12D_N(t)^2 .                 \tag{11}
\]
Это точное тождество в норме Гильберта--Шмидта для двух унитарных операторов. В частности, при $k=59$ нормированный след по модулю больше $0.957$, при $k=57,58$ он больше $0.99999$.

Для исходного веса рациональная проверка даёт $w(3/7)>0.06854$, $w(6/7)>0.01599$, $C_w<0.02067$, $w(0)<0.1141$. Для любого целого $r\ge1$ телескопирование степеней унитарных операторов даёт $\|U_N(t)^{rm}-M_N^{rm}\|_{\rm n}\le r D_N(t)$. Следовательно, при $r=2$ нормированный след на времени $6N/7$ имеет модуль не меньше $1-2D_N(t)^2$. Сохраняя в полном парном тесте оба времени, получаем
\[
 \Delta_{2,N}(b_*)\ge
 w(3/7)(1-D_N(t)^2/2)^2+
 w(6/7)(1-2D_N(t)^2)^2-C_w-\frac{2w(0)}{N}.
\]
При $k\le56$ подставляем $D_N(t)<0.204$, $N\ge49$. При трёх новых $k$ берём $D_N(t)<0.291$, $N\ge7^{57}$. В обоих случаях граница больше $0.053$.

Для ограниченного локального функционала берём $f(y)=(1-|y|/70)_+$, $X=\Xi_N(f)/70$, $Y=\Xi_N(e^{2\pi i (3/7)y}f)/70$, $G=\min(1,X-|Y|)$. Как доказано в предшествующей части, $\EE G\le1-|\Tr U_N^m|/N$ по коэффициенту Фурье от равномерного начала развёртки. У $\si$ имеем $\EE G\ge1-2/\sqrt{105}$ по дисперсии двух линейных статистик. Поэтому на всём диапазоне $2\le k\le59$ доказан разрыв ожиданий больше $0.76$.

\subt{Убывающая сила} Если $t_N=o(N^{-3/5})$, берём $L$ порядка $\min(|t_N|^{-4/3},N/k)$ с малой постоянной, чтобы применить (9). Если минимум равен степени $|t_N|$, правая часть (9) есть $O(N|t_N|^{5/3}+|t_N|^{1/3})=o(1)$. Если минимум задаёт $N/k$, то $|t_N|\le C(k/N)^{3/4}$ и (9) даёт $O(|t_N|\sqrt{Nk})=o(1)$. Округление $L$ в целое не меняет порядков. Тогда (11) даёт нормированный след с модулем, стремящимся к 1, при $u_N=t_N/N=o(N^{-8/5})$. Тот же $G$ исключает $\si$ для этого режима.
\newpage
\sect{6. Спектральные ветви и число переходов между узлами}
Для каждого фиксированного $N$ фазы $U_N(s)$ можно нумеровать непрерывными ветвями $\theta_j(s)$ с локальным аналитическим выбором. Действительно, локально выбирается эрмитов логарифм через дугу, отсутствующую в спектре. Теорема Реллиха о вещественно-аналитических эрмитовых семействах даёт аналитические фазы и векторы там. На компактном интервале соединяем конечное число таких локальных параметризаций. Дифференцирование собственных уравнений при $\dot U_N=iA_N U_N$ даёт $2\pi\dot\theta_j=\langle v_j,A_Nv_j\rangle$. Следовательно, все поднятые фазы липшицевы со скоростью $\le \|A_N\|_{\rm op}/(2\pi)$, включая выбор при совпадающих фазах.

Пусть $M_N^m=e^{2\pi i m\gamma_N}I$. Ветвь начинает движение в узле $\gamma_N+j/m$. Обозначим
\[
  d_i(s)=\operatorname{dist}\!\left(m(\theta_i(s)-\gamma_N),\Z\right),
                 \qquad V_*=\sup_{|s|\le |t_N|}D_N(s).
\]
Из спектральной записи $D_N(s)^2=N^{-1}\sum_i|e^{2\pi im(\theta_i(s)-\gamma_N)}-1|^2$ следуют две оценки
\[
       \frac1N\sum_i d_i(s)^2\le D_N(s)^2/16,\qquad
       \frac{\#\{i:d_i(s)\ge1/4\}}{N}\le D_N(s)^2/2.                \tag{12}
\]
Первая использует $\sin(\pi x)\ge2x$ на $[0,1/2]$. Во второй на $[1/4,1/2]$ квадрат модуля не меньше 2.

Округлим конечные фазы при $s=t_N$ до ближайшего узла той же сетки. Если его номер отличается от номера начального узла, ветвь пересекла середину какого-либо промежутка. Перед первым таким пересечением она провела в области $d_i(s)\ge1/4$ не менее $c/N$ времени: расстояние $d_i$ прошло от $1/4$ до $1/2$ со скоростью $\le C N$. Можно взять последний момент $d_i=1/4$ перед первым моментом $d_i=1/2$. Повторные пересечения считать не требуется. Суммируя длительности по ветвям и применяя (12) в каждый момент, получаем для числа $R_N$ ветвей, сменивших узел,
\[
                           \frac{R_N}{N}
                \le C N|t_N|\, V_*^2 .                            \tag{13}
\]
Знак $t_N$ не меняет оценок. Случай попадания точно в середину также покрыт, при любом выборе ближайшего узла.

Для $t_N=o(N^{-9/13})$ правая часть (13) стремится к нулю. Действительно, при $L\asymp |t_N|^{-4/3}$ из (9) можно оценить $V_*\le C N |t_N|^{5/3}$: добавочный член $|t_N|^{1/3}$ поглощается, так как $N|t_N|^{4/3}\gtrsim k$ в этом случае. Тогда (13) ограничено $C N^3|t_N|^{13/3}$. При $L\asymp N/k$ имеем $|t_N|\lesssim(k/N)^{3/4}$ и $V_*\le C |t_N|\sqrt{Nk}$. Правая часть (13) не больше $C N^2k|t_N|^3=O(N^{-1/4}k^{13/4})$. Во всех случаях $V_*\to0$.
\newpage
\sect{7. Локальный предел и контроль кратностей}
Чистая линейная карта имеет сетку из $m=3N/7$ узлов с кратностями
\[
                     n_j^0=2\sum_{r=0}^{k-1}\mathbf1_{7^r\mid j}
                                      +\mathbf1_{j=0}.
\]
В предшествующей части это выведено диагонализацией $S$ над $\Z/7^k\Z$ и сопряжением $M$ с перестановкой $x\mapsto\gamma x$. В локальной развёртке у неё случайно сдвинутая решётка шага $7/3$. Кратности совместно сходятся к $2(1+\nu_7(B+j))$ по равномерному начальному индексу по модулю $7^r$ при каждом фиксированном $r$, затем $r\to\infty$.

Пусть $n_j$ --- кратности после округления возмущённых фаз. Тогда $\sum_j|n_j-n_j^0|\le2R_N$, так как каждая сменившая узел ветвь может дать две единицы в этой норме. При общем равномерном начале развёртки ожидаемая разность любого локального теста $g\in C_c^1$ между двумя сеточными процессами не выше $2\|g\|_1 R_N/N$. Снятие округления стоит не более
\[
             \|g'\|_1 \sum_i d_i(t_N)/m
                  \le \tfrac{N}{4m}\|g'\|_1 V_*
\]
по интегрированию по началу и (12). При $t_N=o(N^{-9/13})$ обе ошибки стремятся к нулю. Через счётное определяющее семейство тестов и равномерный первый момент на компактах получаем полный локальный предел той же решётки. Физическая сила при этом $u_N=o(N^{-22/13})$.

\subt{Вторые моменты в более узком режиме} В чистой сетке число уровней в любом циклическом интервале из $b$ узлов не больше $C(b+k)$, поскольку среди $b$ индексов кратных $7^r$ не более $b/7^r+1$. По липшицевой границе фаз при движении длины $|t_N|$ в любой конечный округлённый узел могут попасть только ветви из $C(N|t_N|+1)$ соседних исходных узлов. Поэтому
\[
   \max_j(n_j+n_j^0)\le C(N|t_N|+k),\qquad
   \frac1N\sum_j |n_j-n_j^0|^2
                     \le C(N|t_N|+k)\,N|t_N|\,V_*^2 .              \tag{14}
\]
Применим (10) при $L=\lfloor N/k^{1/5}\rfloor$ для достаточно больших $k$, когда $L\le m$. Получим
\[
 V_*\le C\left[
 |t_N|\sqrt N k^{1/10}+t_N^2 N k^{2/5}
                         +|t_N|^3 N^2 k^{1/5}\right].
\]
При $|t_N|=o(N^{-3/4}k^{-1/20})$ правая часть (14) стремится к нулю. Для проверки положим $b_N=N^{3/4}|t_N|$. Множитель $(N|t_N|)^2$ перед $V_*^2$ даёт три верхние границы $C b_N^4 k^{1/5}$, $C N^{-1/2} b_N^6 k^{4/5}$ и $C b_N^8 k^{2/5}$. Все исчезают. Член с $kN|t_N|$ исчезает также, например подстановкой верхней границы $|t_N|\le N^{-3/4}k^{-1/20}$ во все неотрицательные выражения. Таким образом, квадраты разностей кратностей контролируются вплоть до физического режима $u_N=o(N^{-7/4}(\log N)^{-1/20})$.
\newpage
\sect{8. Полный парный предел и воспроизведение}
Завершим переход от оценки кратностей к исходному парному тесту. На любом фиксированном локальном отрезке число точек истинного процесса не превосходит числа округлённых точек на отрезке, расширенном на половину шага сетки с каждой стороны. Расширенное окно охватывает ограниченное число узлов равномерно по $N$. Сумма чистых кратностей на этих узлах имеет равномерно ограниченный четвёртый момент: случайный индекс начала равномерен, а для каждого чистого узла вероятность потребовать больше $r$ разрядов делимости не больше $7^{-r}$ с точностью до постоянного сдвига $r$. Дополнительная неподвижная точка не нарушает оценку. Квадрат суммы абсолютных разностей кратностей на узлах расширенного окна имеет ожидание $o(1)$ по (14) и Коши. Поэтому квадраты локальных счётчиков для исходного и чистого процессов равномерно интегрируемы в указанном режиме. Вместе со слабой сходимостью это даёт сходимость локальных ожиданий парных функционалов.

Для полноты введём стационарную меру разностей $\gamma_N$, включая диагональ. Для $g_b(y)=\int b(x)b(x+y)dx$ имеем $2\Delta_{2,N}(b)=(\gamma_N-\EE_{\rm CUE}\gamma_N)(g_b)$, как выведено в части о возврате Пелля. На компактах сходимость $\gamma_N$ следует из предыдущего абзаца, применяя к стационарному процессу непрерывную функцию двух переменных $\rho(x)g(y-x)$ с $\rho\ge0$, $\int\rho=1$ и компактной опорой. Диагональные совпадения учитываются в обоих процессах. Для хвостов берём $h(x)=(1-|x|)_+$. Стационарность и Коши дают
\[
             (h*\gamma_N)(x)
                  =\EE\,\Xi_N([0,1])\Xi_N([x,x+1])
                  \le \EE\,\Xi_N([0,1])^2=O(1).
\]
Значит, массы всех единичных интервалов ограничены равномерно. Для $b_*$ функция $g_{b_*}$ быстро убывает, поэтому предел сохраняется после снятия компактного усечения. Таким образом, при $u_N=o(N^{-7/4}(\log N)^{-1/20})$
\[
 \lim_{N=7^k\to\infty}\Delta_{2,N}(b_*)=
     \sum_{a=0}^{\infty}7^{-2a}\!
          \sum_{\substack{u\ge1,\ 7\nmid u\\3u<2\cdot7^{a+1}}}
             w(3u/7^{a+1})-\int_0^2\min(t,1)w(t)\,dt .
\]
Правая часть есть точно вычисленный ранее предел для чистой карты. В более широкой области локальной сходимости $u_N=o(N^{-22/13})$ сохраняется нижняя предельная граница той же константой: компактные парные функционалы неотрицательны, к ним применяется усечение сверху, затем монотонное снятие усечений. Равенство там не утверждается.

\textbf{Воспроизведение.} Архив содержит точные проверки группировки коэффициентов по $Q$, взаимной делимости компонент, счётчиков столкновений, трёх рациональных подстановок в (9) и границ наблюдаемых. Параметры сетки и значения $w$ берутся из пройденной проверки обычного следа.

\textbf{Источник теоремы о спектральных ветвях.} A. Rainer, \emph{Perturbation theory of polynomials and linear operators}, теорема 2.1 (Реллих). \href{https://www.mat.univie.ac.at/~armin/publ/handbook_paper.pdf}{Текст автора}.
\newpage
\sect{9. Составные размерности с тем же возвратом}
Оценка (9) также применима на некоторых размерностях $N=d\,7^k$, где $d$ нечётно и взаимно просто с 7. В этом разделе обозначим $N_7=7^k$ и оставим $m=3N_7/7$. Предположим, что $S^m=I\bmod d$. Тогда $S^m=I\bmod N$ по китайской теореме, и $M_N^m$ скалярна. При этом в ограничении $L\le N_7/(1000k)$ используется только компонента $N_7$, а в правой части (9) подставляется $t=N/10^{80}$ с полной размерностью $N$.

Обоснуем перенос. Метки, совпадающие по модулю полной размерности, обязательно совпадают по модулю $N_7$. Поэтому полное совпадение сумм двух меток задаёт уточнение разбиения пар, использованного в (5) для $N_7$. Сумма квадратов мощностей при уточнении может лишь уменьшиться. То же верно для совпадений одиночных меток в (1): разные классы $Q$ по модулю $N_7$ не пересекаются и пара $(0,\pm7)$ имеет разные орбиты уже после редукции. Базис Вейля полной матрицы всё ещё ортонормирован, поэтому нормы в (1),(7) ограничены теми же числами. Интегрирование по частям действует без изменений.

Положим $h_j+v_j\sqrt2=(1+\sqrt2)^j$. Для нечётного $j$ верно $S^j-I=2h_j\bigl(\begin{smallmatrix}h_j&2v_j\\ v_j&h_j\end{smallmatrix}\bigr)$, как в доказательстве пеллева возврата. Возьмём два точных примера:
\[
\begin{array}{c|c|c|c|c}
j&d&k&\sqrt L&\text{верхняя граница для }D_N(t)\\\hline
21&h_{21}/7^2=1114457&54&47\cdot10^{17}&0.019\\
147&h_{147}/7^3&19&774000&0.130 .
\end{array}
\]
В обоих случаях $d$ взаимно просто с 7 и $j\mid m$, так что $S^m=I\bmod d$. Все столбцы проверяются целочисленной рекурсией Пелля и рациональной подстановкой в (9).

В первом примере $N>4\cdot10^{51}$, $m>10^{45}$ и $|\Tr U_N^m|/N>0.9998$. Во втором $N>3\cdot10^{69}$, $m=4885240793731347$ и $|\Tr U_N^m|/N>0.9915$. Время как доля полной размерности здесь равно $3/(7d)$, то есть отличается от $3/7$ в основном семействе. Для исходного веса на узле $1/5488$ точная нижняя граница $w(1/5488)>0.1138$. На первом модуле сохраним первые 50 времён $rm$ с $1\le r\le50$, на втором первые десять. Все эти времена имеют $rm/N\le1/5488$, откуда $w(rm/N)>0.1138$. По унитарному телескопированию и скалярности линейной степени каждый нормированный след не меньше $1-(rD_N(t))^2/2$ по модулю, пока правая часть положительна. Подставляя указанные верхние границы для $D_N(t)$, получаем
\[
     \Delta_{2,N}(b_*)>4.149\quad\text{на первом модуле},\qquad
     \Delta_{2,N}(b_*)>0.582\quad\text{на втором}.
\]
Вычитаемая эталонная сумма ограничена $C_w+2w(0)/N$, как в чистом семействе степеней 7. Все знаменатели, интервальные оценки и конечные суммы проверяются точно. Это оценки для исходной силы с указанными временами возврата, которые составляют разные доли полной размерности.

\newpage
\sect{10. Локальный предел на смешанных модулях}
Зафиксируем нечётное $d$, взаимно простое с 7, для которого $S^{3\cdot 7^b}=I\bmod d$ при некотором целом $b\ge0$. Рассмотрим $N=d\,7^k$, $k\to\infty$, $N_7=7^k$. При $k>b$ имеем $S^{3N_7/7}=I\bmod N$. Аргумент уточнения из предыдущего раздела переносит и оценку (10), если заменить $N$ в правой части оценки столкновений на $N_7$. Поэтому границы для $D_N(t)$ остаются верны с константами, зависящими от фиксированного $d$. В частности, перенос уровней контролируется (12),(13) на сетке $m=3N_7/7$, шаг которой в единицах развёртки равен $7d/3$.

Для вычисления чистых кратностей воспользуемся китайским разложением $x\mapsto(x\bmod N_7,x\bmod d)$. В соответствующем тензорном базисе
\[
 T_N(a,c)=T_{N_7}(a,\bar d\,c)\otimes T_d(a,\overline{N_7}\,c).
\]
На каждом факторе действие на метках сопряжено с $S$ масштабированием второй координаты. Оно симплектично и потому имеет унитарного реализатора, по аргументу с матричными единицами из предшествующей части. Произведение двух реализаторов реализует исходное действие на всех метках, откуда $M_N$ отличается от него лишь общей фазой. На факторе $N_7$ после диагонализации модульной матрицы получаем ту же масштабную перестановку и те же кратности $n_j^0$ из раздела 7. Фазы второго фактора после общего поворота лежат на сетке знаменателя $3\cdot7^b$. На полной сетке длины $m=3\cdot7^{k-1}$ каждая из них сдвигает индексы на некоторое кратное $7^{k-1-b}$.

Для любого фиксированного числа проверяемых разрядов такие сдвиги исчезают при больших $k$. Поэтому чистый локальный закон есть случайно сдвинутая решётка шага $7d/3$ с кратностями $2d(1+\nu_7(B+j))$, где $B$ равномерен в $\Z_7$. Для строгого перехода от усечённых кратностей к полным используем объединение событий: вероятность, что в заданном конечном числе узлов какой-либо из $d$ сдвигов требует более $R$ разрядов, не больше $C_d 7^{-R}$. Каждый из $d$ сдвигов также сохраняет равномерно ограниченные моменты чистых кратностей. Чистое число уровней на $a$ соседних узлах не больше $C_d(a+k)$. Это именно две оценки, нужные для переноса из разделов 7 и 8. Значит, полный локальный закон сохраняется при $u_N=o(N^{-22/13})$, а контроль второго момента при $u_N=o(N^{-7/4}(\log N)^{-1/20})$.

В последней области предел полного парного теста на смешанном модуле равен
\[
  \sum_{\substack{u\ge1\\3u<14d}} w(3u/(7d))
   +\sum_{a=1}^\infty 7^{-2a}
      \sum_{\substack{u\ge1,\ 7\nmid u\\3u<2d\,7^{a+1}}}
         w(3u/(d\,7^{a+1}))
    - C_w .
\]
Действительно, для чистых следов размер ядра $S^m-I$ перемножается по двум компонентам. В каждом фиксированном классе $a$ второй множитель равен $d^2$ для всех достаточно больших $k$, а после нормировки на $N^2$ первый даёт $7^{-2a}$. При $a=0$ подходят все скалярные времена. Хвост по $a$ ограничен через $C_d\sum_{a\ge A}7^{-a}$, так как ядро на второй компоненте никогда не больше $d^2$. Этого достаточно для снятия ограничения на число классов. Уже первый член суммы превышает $C_w$, поскольку $w(3/(7d))\ge w(3/7)>C_w$. Так получаем семейство различных арифметических локальных законов для вспомогательных убывающих киков. Неизменная сила в (9) использована для конечных примеров предыдущего раздела.
\end{document}
